\subsection{Study 1: Determining the Amount of Image Resources Required for Generalization}

The \textit{Calibration} subset was used to "teach" the CV system how the ant morphology appears in images, while the other subsets were employed to evaluate its generalization capabilities. The “A” subsets were designed to mimic the system being deployed in an ongoing and repetitive experiment, with images with similar imaging backgrounds over time. In contrast, the “B” subsets aim to test the system robustness in handling images with less controlled imaging conditions. This study focused on the performance of the CV system with different numbers of available images for the model calibration. The available images were randomly sampled from the \textit{Calibration} subset given the numbers of 64, 256, and 1024. Since the number of available images was only 954, when the sampled number was set as 1024, the system sampled the images with replacement until the number of images reached 1024. Since the data augmentation that was applied to the images to increase the image diversity, even two images sampled from the same image can be different. The sampling process was repeated 50 times for each subset and each available number to avoid sampling bias. In addition to reporting the model detection metrics that are cmoonly used in machine learning but less intepretable, this study also reported the correlation between the automated and manual counts. The automated counts were obtained by summing the number of ants detected with confidence scores above 0.5, while the manual counts were derived from teh manual YOLO annotations of the images. The correlation was calculated using the Pearson correlation coefficient, which measures the linear relationship between the two sets of counts. A high correlation indicates that the automated counting results are consistent with the manual counts, while a low correlation suggests discrepancies between the two methods.

Additionally, when objects of interest are small, previous studies have demonstrated that model performance can be improved using Slicing Aided Hyper Inference (SAHI) \cite{akyon_slicing_2022}. This technique divides input images into smaller patches, enabling the model to focus on localized areas and potentially improve detection accuracy for small objects. In this study, we evaluated model performance both with and without SAHI on the test images. SAHI implementation requires two key parameters: patch size and overlap size. The patch size determines the dimensions of each image patch, while overlap size controls the degree of overlap between adjacent patches. We explored three patch sizes by setting dividers of 2, 3, and 4, where the patch size equals the shorter side of the image divided by the chosen divider. For example, when the divider is set to 2, the patch size becomes half the length of the image's shorter side. Similarly, we tested three overlap sizes: 0, 0.25, and 0.5, where an overlap size of 0 indicates no overlap between patches.

Since each combination of these parameters potentially yields different performance outcomes, we conducted a grid search by evaluating all nine possible combinations of patch size and overlap size through hyperparameter tuning. Before testing the model on each subset, we sampled 10 images from the target subset to determine the optimal parameter combination for SAHI implementation. Given the absence of ground truth data for ant counts in these tuning images, we used high-confidence detections (confidence > 0.5) as a proxy for prediction quality. This unsupervised approach avoids data leakage issues since ground truth labels are not used during the hyperparameter tuning process.

As the investigated factors, including the calibration sample size, with and without SAHI, model complexity, and the background similarity may affect the model generalization performance, an ordinary least squares (OLS) regression model was used to evaluate the impact of these factors on the model performance:

\begin{equation}
y_{ijk} = \beta_0 + \beta_1 n_i + \beta_2 m_j + \beta_3 b_k + \varepsilon_{ijk}
\label{eq4}
\end{equation}

where $y_{ijk}$ is the model performance metric (e.g., mAP@0.5) for the $i$-th calibration sample size, $j$-th model complexity, and $k$-th background similarity, $\beta_0$ is the intercept, $\beta_1$, $\beta_2$, and $\beta_3$ are the coefficients for calibration sample size, model complexity, and background similarity, respectively, and $\varepsilon_{ijk}$ is the error term. The calibration sample size $n_i$ is a categorical variable with three levels: 64, 256, and 1024; the model complexity $m_j$ is a categorical variable with three levels: YOLOv11n, YOLOv11m, and RT-DETR-L; and the background similarity $b_k$ is a categorical variable with two prefixes: "A" and "B". The regression model was fitted using the Python Statsmodels library \cite{seabold_statsmodels_2010}.

Additionally, the performance with SAHI and without SAHI was evaluated by statistically comparing their mean Average Precision (mAP) metric over the 50 iterations. We began by examining data distribution characteristics using the Shapiro-Wilk normality test \cite{shapiro_analysis_1965} for each performance distribution. When both distribution demonstrated normal distributions, we subsequently assessed variance homogeneity through Levene's test \cite{levene_robust_1960}.  The choice of parametric test depended on these preliminary assessments: Student's t-test was applied when assumptions of normality and equal variances were satisfied, while Welch's unequal variances t-test \cite{welch_generalization_1947} was used when variances differed significantly despite normal distributions. For performance distribution that failed normality assumptions, we employed the non-parametric Mann-Whitney U test \cite{mann_test_1947} to evaluate distributional differences. All statistical inferences were made using $\alpha$ = 0.05 as the threshold for significance. The statistical analysis was conducted using the Python SciPy library \cite{virtanen_scipy_2020}.
